\documentclass[11pt]{article}
\usepackage{mathblatt}
\begin{document}
\weit
\typeout{PROBE-BEGIN winkel-dreiecke-e2-k1-s0-v1}
\begin{aufgabe}{\detokenize{winkel-dreiecke-e2-k1-s0-v1}}
\begin{teile}
\teil Probe

\parallelenpaar{58}{}{}{}{}
\end{teile}
\end{aufgabe}
\typeout{PROBE-END winkel-dreiecke-e2-k1-s0-v1}
\typeout{PROBE-BEGIN winkel-dreiecke-e2-k1-s0-v2}
\begin{aufgabe}{\detokenize{winkel-dreiecke-e2-k1-s0-v2}}
\begin{teile}
\teil Probe

\parallelenpaar{67}{}{}{}{}
\end{teile}
\end{aufgabe}
\typeout{PROBE-END winkel-dreiecke-e2-k1-s0-v2}
\typeout{PROBE-BEGIN winkel-dreiecke-e2-k1-s0-v3}
\begin{aufgabe}{\detokenize{winkel-dreiecke-e2-k1-s0-v3}}
\begin{teile}
\teil Probe

\parallelenpaar{71}{}{}{}{}
\end{teile}
\end{aufgabe}
\typeout{PROBE-END winkel-dreiecke-e2-k1-s0-v3}
\typeout{PROBE-BEGIN winkel-dreiecke-e2-k1-s0-v4}
\begin{aufgabe}{\detokenize{winkel-dreiecke-e2-k1-s0-v4}}
\begin{teile}
\teil Probe

\parallelenpaar{52}{}{}{}{}
\end{teile}
\end{aufgabe}
\typeout{PROBE-END winkel-dreiecke-e2-k1-s0-v4}
\typeout{PROBE-BEGIN winkel-dreiecke-e2-k2-s0-v3}
\begin{aufgabe}{\detokenize{winkel-dreiecke-e2-k2-s0-v3}}
\begin{teile}
\teil Probe

\parallelenpaar{62}{}{}{}{}
\end{teile}
\end{aufgabe}
\typeout{PROBE-END winkel-dreiecke-e2-k2-s0-v3}
\end{document}
